\labsection{1}{Processes and the Linux namespace}{sec:lab01}
\begin{labproject}{From one executable to observable processes and directory names}

\textbf{Coverage.} Chapter~\ref{ch:processes}: PID/PPID, \texttt{fork()}, directory names.\\
\textbf{Goal.} Run the exact supplied source and inspect live OS state.\\
\textbf{Explain.} Explain how one executable becomes two processes and how directory names change.

\courselabflow{Build programs}{Inspect directory}{Fork process}{Wait and verify}



\section*{Source code: read before you run}
\labsourcecard{\texttt{directory\_list.c}}{Open a directory stream and print every entry returned by Linux.}{The acquire/use/release sequence: \texttt{opendir} $\rightarrow$ \texttt{readdir} loop $\rightarrow$ \texttt{closedir}.}
\noindent\textbf{Start the file.} \texttt{directory\_list.c}, lines 1--33.
\courselabtrace{One fork, two return paths}{Parent}{fork returns the positive child PID}{Child}{fork returns zero; getppid identifies its parent}{Join}{waitpid collects the child exit status}

\lstinputlisting[style=utpcnumbered,firstline=1,lastline=33]{all_experiments/Lab01/directory_list.c}

\noindent\textbf{Continue the same file.} \texttt{directory\_list.c}, lines 34--35.
\lstinputlisting[style=utpcnumbered,firstline=34,lastline=35]{all_experiments/Lab01/directory_list.c}

\labsourcecard{\texttt{fork\_demo.c}}{Create a child process, print process identity, optionally keep the child alive, then wait for it.}{The return value from \texttt{fork()}, the child's PPID, and the parent's \texttt{waitpid}.}
\noindent\textbf{Start the file.} \texttt{fork\_demo.c}, lines 1--34.


\textbf{Use the chapter while coding.} Revisit \hyperref[ch:processes]{the earlier chapter explanation}; match its operation to the variables and state changes below.

\textbf{Follow the supplied implementation.} Read the branch condition before attributing a line of output to parent or child. Scheduling can change print order. Use the optional delay for observation, then verify the parent waits rather than inferring behaviour from a single output ordering.
\courselabtrace{Trace fork\_demo and process ownership}{1. Input / state}{fork returns a child PID in the parent and zero in the child; failure is a separate branch.}{2. Operation}{The child prints its identity; the parent retains child\_pid and waits for that specific child.}{3. Check}{waitpid collects the child's termination status; printed PID values vary by run.}


\lstinputlisting[style=utpcnumbered,firstline=1,lastline=34]{all_experiments/Lab01/fork_demo.c}

\noindent\textbf{Continue the same file.} \texttt{fork\_demo.c}, lines 35--51.
\lstinputlisting[style=utpcnumbered,firstline=35,lastline=51]{all_experiments/Lab01/fork_demo.c}

\begin{utptable}
\renewcommand{\arraystretch}{1.18}
\begin{tabularx}{\linewidth}{@{}>{\bfseries\leavevmode\color{UTPNavy}}m{0.34\linewidth}
  >{\raggedright\arraybackslash}X@{}}
\toprule
\rowcolor{UTPPaleBlue!92}
\utpth{Code region} & \utpth{What it means} \\
\midrule
\texttt{opendir} / \texttt{readdir} / \texttt{closedir} & Linux exposes directory entries through a stream-like API. \\
\texttt{child\_pid = fork()} & One running process becomes parent + child. \\
\texttt{child\_pid == 0} & This branch is executing in the child. \\
\texttt{waitpid(...)} & The parent waits and collects the child's exit status. \\
\bottomrule
\end{tabularx}
\end{utptable}

\section*{Task 1 --- Build a clean baseline}
\begin{lstlisting}[style=utpbash]
cd all_experiments/Lab01
gcc -Wall -Wextra -std=c11 directory_list.c -o directory_list
gcc -Wall -Wextra -std=c11 fork_demo.c -o fork_demo
./directory_list .
./fork_demo
\end{lstlisting}
Do not change the source yet. Record the parent PID, child PID, and child PPID from one run.

\section*{Task 2 --- Observe the namespace changing}
Run \texttt{./directory\_list .} after each command:
\begin{lstlisting}[style=utpbash]
touch observation.txt
mv observation.txt renamed.txt
rm renamed.txt
\end{lstlisting}
Create a four-row table: initial state, after \texttt{touch}, after \texttt{mv}, after \texttt{rm}. Record only the name that changed.

\section*{Task 3 --- Observe the parent and child while alive}
\begin{lstlisting}[style=utpbash]
./fork_demo 10 &
ps -C fork_demo -o pid,ppid,state,comm
wait
\end{lstlisting}
Match the child's PPID to the parent's PID. The state letter can vary because scheduling is real and timing-dependent.

\section*{Task 4 --- Controlled edit in Vim}
Open \texttt{fork\_demo.c} with Vim. Change the child message so it includes the words \texttt{child branch observed}. Save, recompile, and run again.

Before running, answer: \emph{Will the parent print the new child message too?} Then explain the observed result using the idea of two processes, not only the \texttt{if} statement.

\begin{quickcheck}
\begin{enumerate}
  \item Why does \texttt{fork()} produce two different return-value cases from one call?
  \item How can a child process identify its parent?
  \item Why can creating or renaming a file change a directory listing without changing the directory program itself?
\end{enumerate}
\end{quickcheck}

\textbf{Chapter connection.} Revisit \hyperref[ch:processes]{the chapter explanation} and map its states to the steps below.
% Reused from chapters/02-linux-processes-files-and-system-calls.tex
\begin{labvisual}{Recall the chapter model}
\begin{labgraphic}
\begin{tikzpicture}[
  box/.style={draw=UTPBlue,fill=UTPPaleBlue,rounded corners=2mm,minimum width=3.5cm,minimum height=1cm,align=center,font=\small\color{UTPNavy}},
  live/.style={box,draw=UTPTeal,fill=UTPPaleTeal},
  kern/.style={box,draw=UTPNavy,fill=UTPNavy!10},
  ar/.style={-{Stealth[length=2.2mm]},thick,draw=UTPOrange}]
\node[box] (file) at (-4.2,0) {Executable file\\\texttt{fork\_demo}};
\node[live] (proc) at (0,0) {Running process\\PID 4200};
\node[kern] (os) at (4.2,0) {Kernel process table\\state, PID, PPID};
\draw[ar] (file) -- node[above,font=\scriptsize]{run} (proc);
\draw[ar] (proc) -- (os);
\end{tikzpicture}
\end{labgraphic}\end{labvisual}


\begin{verification}
\textbf{Explain and extend your result.} Start a harmless long-running command, record its PID and parent, terminate only that process and confirm its exit. Explain why a second process with the same command name must remain unaffected.


\begin{itemize}
  \item Both programs compile with \texttt{-Wall -Wextra} and run from the unchanged source first.
  \item The child's PPID matches the parent's PID in the \texttt{ps} output.
  \item The namespace table records exactly one changed name after each of \texttt{touch}, \texttt{mv} and \texttt{rm}.
  \item Your prediction for the Vim edit is written down before you run it, and the result is explained with two processes.
\end{itemize}
\end{verification}

\begin{troubleshoot}
If \texttt{ps -C} shows nothing, rerun \texttt{./fork\_demo 10 \&} and issue \texttt{ps} before the child exits. Different PID numbers are expected on every run.
\end{troubleshoot}
\end{labproject}
